% The physics of the symmetry data: representations, fusion, invariance, and the container.
\chapter{Tensors}
We follow the convention by \cite{tenpy2024}.
\section{Symmetry and why representations appear}

A quantum state of a lattice model is a vector in a Hilbert space
$\mathcal H=\bigotimes_{i}\mathcal H_i$, and a symmetry group $G$ acts
on it through a representation $U(g)=\bigotimes_i u_i(g)$.
The sites'' local spaces $\mathcal H_i$ are therefore representations of $G$,
and the physically meaningful operators are those that transform in a
definite way under $G$.

For a compact Lie group the irreducible representations (irreps) are the
elementary building blocks.  Any local space decomposes as
\begin{equation}
  \mathcal H_i \;\cong\; \bigoplus_{\lambda} m^{i}_{\lambda}\, V_{\lambda},
\end{equation}
with $V_{\lambda}$ the irrep of highest weight $\lambda$ and $m^i_\lambda$ its
multiplicity.  In a spin chain, for example, the local space is a single irrep:
$\mathcal H_i = V_j$ for spin $j$.  In a Hubbard or $t$--$J$ model it is a
direct sum of several irreps.

Two facts organize everything that follows.

\begin{enumerate}
  \item \emph{Only the multiplicities are data.}  An irreducible sector
        $V_\lambda$ of a local space has no free parameters: once the group
        and the irrep are fixed, the matrix elements of the generators are the
        universal representation matrices.
  \item \emph{Intertwiners carry the physics.}  A tensor that is invariant
        under $G$ is a map between tensor products of irreps that commutes
        with the action, and by Schur''s lemma such a map is a sum of
        proportional copies of the few elementary intertwiners the fusion
        rules allow.  The proportionality constants are the physical numbers.
\end{enumerate}

The rest of these notes makes those two statements precise, then describes
the block-sparse container that realizes them, and the cost argument for why
the sparse route is the production algorithm.

\section{Lie algebras and their representations}
\label{sec:lie}

\subsection{Generators, structure constants, Casimir}

A Lie group $G$ has a Lie algebra $\mathfrak g$ with generators $T^a$,
$a=1,\dots,\dim\mathfrak g$, satisfying
\begin{equation}
  [T^a,T^b] \;=\; \mathrm i\, f^{ab}{}_{c}\, T^{c},
\end{equation}
with the totally antisymmetric structure constants $f^{abc}$.  The quadratic
Casimir
\begin{equation}
  C_2 \;=\; \sum_{a} T^a T^a
\end{equation}
commutes with every generator, $[C_2,T^a]=0$, so by Schur''s lemma it is a
constant on each irrep.  It is the natural label-independent measure of an
irrep''s ``size''.

A maximal set of commuting generators $H_i$, $i=1,\dots,r$ (the rank of
$\mathfrak g$ being $r$) spans the Cartan subalgebra.  Their simultaneous
eigenvalues are the {\em weights}: a state of weight $\mu$ satisfies
$H_i|\mu\rangle=\mu_i|\mu\rangle$.  Raising and lowering operators
$E_\alpha$, associated with the roots $\alpha$, move between weights,
$\mu\mapsto\mu+\alpha$; positive roots raise, and a highest weight
$\lambda$ is annihilated by all of them.

\subsection{Highest-weight representations and the dimension formula}

An irrep is built from a highest weight: one starts from the state
$|\lambda\rangle$ with $E_{\alpha}|\lambda\rangle=0$ for all positive roots
$\alpha$, and applies lowering operators repeatedly.  The result is a
finite-dimensional irrep iff $\lambda$ is dominant integral, which in
Dynkin labels means non-negative integer coefficients on the fundamental
weights.

The dimension is the Weyl formula
\begin{equation}
  \dim V_\lambda \;=\; \prod_{\alpha>0}
  \frac{\langle \lambda+\rho,\alpha\rangle}{\langle \rho,\alpha\rangle},
  \qquad \rho \;=\; \tfrac12\sum_{\alpha>0}\alpha ,
  \label{eq:weyl}
\end{equation}
and the Casimir eigenvalue is
\begin{equation}
  C_2(\lambda) \;=\; \langle \lambda,\lambda+2\rho\rangle ,
\end{equation}
in the normalization in which the generators are those of the defining
representation.

\subsection{The two groups at hand}

\paragraph{$\SU{2}$.}  Rank one, one Cartan generator $J^z$, and
irreps labelled by the spin $j=0,\tfrac12,1,\dots$, of dimension
\begin{equation}
  \dim V_j = 2j+1, \qquad C_2 = j(j+1).
\end{equation}
Weights run over $m=-j,-j+1,\dots,j$; the lowering operator $J^-$ maps
$|j,m\rangle$ to $|j,m-1\rangle$ with the known matrix element
$\sqrt{(j+m)(j-m+1)}$, and the raising operator is its adjoint.
Internally the package labels a spin by the doubled integer $2j$, which is
an integer and therefore hashable and comparable; the dimension of the
sector is then $2j+1$.

\paragraph{$\SU{3}$.}  Rank two, two commuting Cartan generators $T^3$, $T^8$
(the diagonal Gell-Mann matrices) and irreps labelled by the Dynkin labels
$(p,q)$, $p,q\ge 0$.  The highest weight is
\begin{equation}
  \lambda = p\,\omega_1 + q\,\omega_2,
\end{equation}
with $\omega_{1,2}$ the fundamental weights; the weights of the irrep form a
hexagonal (or triangular) pattern in the $(T^3,T^8)$ plane, and the
fundamental weights are the corners.  The dimension is
\begin{equation}
  \dim V_{(p,q)} = \frac{(p+1)(q+1)(p+q+2)}{2},
\end{equation}
and the Casimir, in the normalization of the Gell-Mann generators
$\operatorname{tr}\lambda_a\lambda_b = 2\delta_{ab}$,
\begin{equation}
  C_2(p,q) = \frac{p^2+q^2+pq+3p+3q}{3},
\end{equation}
so that the defining irrep has $C_2=\tfrac43$ and the adjoint $C_2=3$.
The conjugate of $(p,q)$ is $(q,p)$: quarks and antiquarks, and the
generators in the conjugate basis are $-\lambda^T$ at the level of the
defining representation, but for a general irrep they are the generators of
the {\em independently constructed} conjugate irrep, not the negated
generators of the original one.  That distinction is the one place where a
purely algebraic guess fails, see section~\ref{sec:conj}.

\section{Tensor products, fusion and Clebsch--Gordan}
\label{sec:fusion}

\subsection{Decomposition and multiplicities}

The tensor product of two irreps decomposes as
\begin{equation}
  V_\lambda \otimes V_\mu \;\cong\; \bigoplus_{\nu} N^{\nu}_{\lambda\mu}\, V_{\nu},
  \label{eq:fusion}
\end{equation}
where the fusion multiplicity $N^{\nu}_{\lambda\mu}$ is a non-negative
integer.  Groups with $N^{\nu}_{\lambda\mu}\in\{0,1\}$ for all triples are
called multiplicity-free; $\SU{2}$ and $\SU{3}$ have this property for the
products of their defining representations, but not in general:
\begin{equation}
  \tfrac12\otimes\tfrac12 = 0\oplus 1, \qquad
  \mathbf 3\otimes\mathbf 3 = \mathbf 6 \oplus \bar{\mathbf 3}, \qquad
  \mathbf 8\otimes\mathbf 8 = \mathbf{27}\oplus\mathbf{10}\oplus
  \overline{\mathbf{10}}\oplus 2\cdot\mathbf 8\oplus\mathbf 1 .
  \label{eq:examples}
\end{equation}
The last line shows a genuine multiplicity $2$: there are two independent
invariant ways to couple two adjoints into an adjoint, and a representation
scheme that cannot carry that label cannot represent $\SU{3}$ physics.

\paragraph{The $\SU{2}$ case.}  $V_{j_1}\otimes V_{j_2} = \bigoplus_{j=|j_1-j_2|}
^{j_1+j_2} V_j$ with every multiplicity one; the coupling is
$m=m_1+m_2$.

\paragraph{The $\SU{3}$ case.}  The multiplicities follow from the
Klimyk--Brauer formula,
\begin{equation}
  N^{\nu}_{\lambda\mu} \;=\;
  \sum_{w\in W} (-1)^{\ell(w)}\,
  \mathrm{mult}_{\lambda+\mu}\!\left(\nu+\rho-w\rho\right),
  \label{eq:klimyk}
\end{equation}
a signed sum over the Weyl group $W$ (order six for $\SU{3}$), where
$\mathrm{mult}$ counts weight multiplicities in the product $\lambda+\mu$.
The formula is exact and computationally cheap once the weights are known;
it is what the package evaluates, in integer coordinates on the weight
lattice so that the comparison of dominant weights is exact.

\subsection{Clebsch--Gordan coefficients}

For a fixed triple $(\lambda,\mu,\nu)$ and a fixed copy index $c$ when
$N^{\nu}_{\lambda\mu}>1$, the Clebsch--Gordan (CG) coefficients are the
matrix elements of the projector that selects $V_\nu$ inside
$V_\lambda\otimes V_\mu$.  They are a third-rank tensor
\begin{equation}
  C^{(\lambda\mu\nu)}_{a b}{}^{c}, \qquad a=1,\dots,\dim V_\lambda,\;
  b=1,\dots,\dim V_\mu,\; c=1,\dots,\dim V_\nu ,
\end{equation}
characterized up to normalization by invariance,
\begin{equation}
  \left(T^{(\lambda)}_x\otimes 1 + 1\otimes T^{(\mu)}_x\right) C
  \;=\; C\,T^{(\nu)}_x \quad\text{for every generator } x,
  \label{eq:cg-invariance}
\end{equation}
i.e.\ $C$ is an intertwiner $V_\lambda\otimes V_\mu\to V_\nu$.  The standard
conventions (Condon--Shortley for $\SU{2}$, the Gell-Mann basis for $\SU{3}$)
fix the normalization and a phase; the package normalizes each CG tensor so
that its Frobenius norm squared equals the dimension of the coupled irrep,
\begin{equation}
  \|C^{(\lambda\mu\nu)}\|^2 = \dim V_\nu ,
  \label{eq:cg-norm}
\end{equation}
which makes the orthogonal-group relations below come out with plain
integer factors.  For $\SU{2}$ this is the usual relation between the CG
coefficients and the Wigner $3j$ symbols,
\begin{equation}
  \langle j_1 m_1 j_2 m_2 | j m\rangle
  = (-1)^{j_1-j_2+m}\sqrt{2j+1}\,
  \begin{pmatrix} j_1 & j_2 & j \\ m_1 & m_2 & -m \end{pmatrix}.
\end{equation}

\subsection{Recoupling and $6j$ symbols}

Coupling three irreps in different orders gives bases related by the
recoupling (Racah) coefficients, the $\SU{2}$ $6j$ symbols
\begin{equation}
  \begin{Bmatrix} j_1 & j_2 & j_3 \\ j_4 & j_5 & j_6 \end{Bmatrix},
\end{equation}
and for a general group by the corresponding invariant intertwiners.  They
are the change-of-basis matrices on the space of intertwiners
$V_a\otimes V_b\otimes V_c\to V_d$, and they are what a network needs
whenever a fusion tree is re-associated, e.g.\ when a leg is split or a
tensor is contracted along a different pairing.  The package computes them
for $\SU{2}$ from the CG tables; for $\SU{3}$ the analogous data follow from the
invariance condition~\eqref{eq:cg-invariance}.

\subsection{Conjugate representations and the invariant form}
\label{sec:conj}

Every irrep $\lambda$ has a conjugate $\bar\lambda$ whose weights are the
negatives.  For a self-conjugate irrep ($\lambda=\bar\lambda$, e.g.\ every
integer-spin $\SU{2}$ irrep and the $\SU{3}$ adjoint) there exists an invariant
bilinear form
\begin{equation}
  D_{\lambda}: V_{\lambda}\otimes V_{\lambda}\to \mathbb C,\qquad
  D_{\lambda}(u(g)v,u(g)w) = D_{\lambda}(v,w),
\end{equation}
unique up to scale, the ``metric''.  It is symmetric or antisymmetric, and
the sign is the Frobenius--Schur indicator
\begin{equation}
  \varepsilon_\lambda = (-1)^{2j} = \begin{cases}
    +1, & \text{integer spin (orthogonal)},\\
    -1, & \text{half-integer spin (symplectic)} .
  \end{cases}
  \label{eq:fs}
\end{equation}
Two consequences matter for a network.  First, a leg that is oriented
oppositely (a ``dual'' leg) is related to the primal one by $D_\lambda$;
this is why the contraction of two legs of opposite orientation is the
invariant pairing, not the naive sum over equal indices.  Second, for
non-self-conjugate irreps the two bases are genuinely different vector
spaces, and the intertwiner between $V_\lambda$ and $V_{\bar\lambda}$ is a
non-trivial matrix.  A negative $\varepsilon$ in a real basis corresponds to
the antisymmetric invariant of the symplectic rep; getting this sign wrong
turns an Hermitian effective Hamiltonian into a non-Hermitian one, which is
exactly the kind of error that survives every symmetry-blind test.

\section{Invariant tensors and symmetric networks}
\label{sec:invariant}

\subsection{The trivial representation}

$G$-invariance means transforming in the trivial irrep $\nu=0$.  A scalar
built from local operators is invariant iff its fusion to the trivial
channel is non-zero,
\begin{equation}
  \mathcal O \text{ invariant} \iff 0 \in \lambda_1\otimes\dots\otimes\lambda_k ,
\end{equation}
and the number of independent invariants is the corresponding fusion
multiplicity.  For a Hamiltonian the total charge is trivial by definition,
which is a checkable property of the network and not an assumption about
the local terms.

\subsection{Wigner--Eckart and reduced matrix elements}

The Wigner--Eckart theorem is the reason a sparse representation is not just
possible but natural.  For a tensor operator $T^{(k)}_q$ transforming in the
irrep $k$,
\begin{equation}
  \langle j' m'| T^{(k)}_q | j m\rangle
  = \langle j\, m;\, k\, q | j' m'\rangle\,
  \frac{\langle j'\|T^{(k)}\|j\rangle}{\sqrt{2j'+1}},
  \label{eq:we}
\end{equation}
the entire $m$-dependence is a CG coefficient, and the physics sits in the
{\em reduced} matrix element $\langle j'\|T^{(k)}\|j\rangle$, a single
number per triple of irreps.  In the general case one has
\begin{equation}
  \langle \alpha' \nu c'| T^{(\lambda_1\dots\lambda_k)}_{\text{inv}}| \alpha \mu c\rangle
  \;=\; \sum_{\text{intertwiners}} \text{(CG data)} \times (\text{reduced data}),
\end{equation}
with the reduced data again the only free parameters.

\subsection{Why every tensor of a symmetric network is an intertwiner}

Consider a tensor network whose legs carry irreps and whose local tensors are
invariant.  Each local tensor is then a map
\begin{equation}
  T: V_{\lambda_1}\otimes\dots\otimes V_{\lambda_p}
  \longrightarrow V_{\mu_1}\otimes\dots\otimes V_{\mu_q},
  \qquad T\,U(g) = U(g)\,T ,
\end{equation}
i.e.\ an intertwiner.  By Schur''s lemma it is a sum of the elementary
intertwiners allowed by the fusion rules, each multiplied by a reduced
matrix element.  Applied to a state, the network therefore carries a
definite total charge at every bond: the bond is not a featureless index
but the label of an irrep with a definite position in the fusion chain.

\section{The block-sparse representation}
\label{sec:container}

\subsection{Legs, charges and blocks}

A leg of a tensor is a labelled collection of irrep sectors,
\begin{equation}
  L = \bigl\{ (s_1, d_1), \dots, (s_n, d_n)\bigr\},\qquad d_i = m_i \dim V_{s_i},
\end{equation}
where $s_i$ is the sector (the charge), $m_i$ the multiplicity and $d_i$ the
number of indices.  The bond dimension a state reports is the sum of the
multiplicities, $\sum_i m_i$: this is the $D$ of the QSpace convention, which
counts multiplets, and it is the unit of the truncation bound.  The number of
states on the leg is $\sum_i d_i = \sum_i m_i \dim V_{s_i}$ instead.  A tensor
is a set of blocks, each block labelled by the
sector chosen on every leg, and the block data are the reduced matrix
elements of the corresponding intertwiner.  The label of a block consists of
\begin{itemize}
  \item the sector index on every leg, and
  \item for every internal vertex of the fusion chain, the intermediate
        sector and (when the group has fusion multiplicities) the copy
        index that selects which of the $N^{\nu}_{\lambda\mu}$ intertwiners
        is used.
\end{itemize}
The block shape is the product of the multiplicities only; the CG tensors,
which are universal and shared by every block with the same fusion tree, are
inserted when a block is expanded into indices.  This is the quantitative
content of Wigner--Eckart: the spin-$j$ structure is not stored, it is
looked up.

\subsection{Fusion trees}

The chain of intermediate sectors is recorded by a {\em fusion tree}: an
ordered list of the leg sectors together with the inner couplings and the
total charge,
\begin{equation}
  \text{tree} = \bigl(s_1,\dots,s_k \,\big|\, c_1,\dots,c_{k-2} \,\big|\, c\bigr),
  \label{eq:tree}
\end{equation}
subject to $c_i \in s_1\otimes\dots\otimes s_{i+1}$ and $c\in
s_1\otimes\dots\otimes s_k$.  A tensor whose legs are chosen admits one
block per admissible tree; for a multiplicity-free group the tree is
determined by the sectors, and for a group with multiplicities each copy
index multiplies the number of admissible trees.  The CG tensor of a tree is
built once and cached; expanding it multiplies the vertex CG tensors along
the chain.

\subsection{Cost: why the sparse route is the production path}

The dense Hilbert space of a bond of dimension $D$ has $D^2$ entries;
the same bond for a group with $r$ irreps of dimensions $d_\alpha$ and
multiplicities $m_\alpha$ has $\sum_\alpha m_\alpha^2$ reduced entries,
because Wigner--Eckart factors out the $\sim d_\alpha^2$ entries of each
irrep.  For $\SU{2}$ at spin $j$ the saving is a factor $(2j+1)^2$ per leg;
for a chain of $L$ such legs it is exponential in $L$.  A concrete
measurement of the same physical state in this package gives an 11.5-fold
smaller tensor at $L=16$ with bond dimension $64$, and an 8.5-fold shorter
sweep than the symmetry-blind engine, with the gap growing with the bond
dimension.

Two corollaries are worth stating explicitly.

\begin{itemize}
  \item \emph{The dense expansion is a reference, not an algorithm.}  A
        block-sparse tensor can always be expanded into plain indices by
        re-inserting the CG data.  That expansion is exact and useful for
        testing a network against a dense calculation, but it defeats the
        whole point: it restores the $D^2$ cost and materializes numbers
        that the sparse algorithm never needs.  Expansions therefore belong
        in tests (as an oracle) and in the small-$D$ diagnostic path, never
        in the sweep.
  \item \emph{Every operation is closed under the symmetry.}  Adding tensors
        with equal legs, contracting them, splitting them with an SVD and
        applying them to a state all preserve the block structure, provided
        the same tree conventions are used on both sides of the operation.
        The algorithms are therefore written once, over the sparse labels,
        instead of once per representation.
\end{itemize}

\section{Charge flow through a network}
\label{sec:flow}

\subsection{Bonds carry the running charge}

In a matrix product state the physical leg of each site is an irrep of $G$,
and the bond legs carry the running charge.  A symmetric MPS is a network of
intertwiners whose total charge, read from the fused tree of the whole
chain, is the state''s quantum number:
\begin{equation}
  |\psi\rangle \;=\; \sum_{\{\,q\,\}} \prod_i T^{(i)}_{\text{inv}},
  \qquad q_{\text{tot}} = \text{fusion of all legs} .
\end{equation}
A ground state obtained by a symmetric DMRG run lives in a definite sector:
the trivial one for a Hamiltonian with a singlet ground state, and the
appropriate one otherwise.  This is physics, not bookkeeping: the total
charge cannot be changed by an invariant operator, and an operator that
carries non-trivial charge produces a state in a different sector.

\subsection{Applying an operator}

An operator $\mathcal O$ acting on the physical index transforms in some
irrep $\rho$; the state $\mathcal O|\psi\rangle$ then lives in the sector
$\rho\otimes q_{\text{tot}}$.  In the sparse representation this requires
that the intermediate objects of the application be allowed to carry
non-trivial charge - the bond structure does this automatically once every
tensor''s block labels are derived from the fusion of its own legs, rather
than demanded to be trivial.  A local operator that is a single component of
a multiplet ($S^z$ inside the spin-$1$ triplet, or $\lambda_8$ inside the
$\SU{3}$ adjoint) is exactly such an object: $T^{(k)}_q$ is an intertwiner
$V_j\to V_j\otimes V_k$ whose charge is the irrep $k$, so the multiplet
index is part of the operator''s data.

An operator that carries non-trivial charge is the one place where the two
conventions of the container are not interchangeable, and it is worth stating
why.  Applying an MPO alternates a contraction with a split.  The contraction
of a state tensor with an operator tensor produces an object whose legs can
fuse to the output charge along more than one path, and \emph{which} path is
taken is data: it is what distinguishes $\mathcal O|\psi\rangle$ from any other
state with the same bond dimensions.  A leg in the resolved convention stores
the magnetic index and therefore sums the paths into one dense window
(section~\ref{sec:flow-contract}); a leg in the reduced convention stores one
entry per path (section~\ref{sec:flow-project}).  The reduced convention is
therefore the production route for a charged operator, and the resolved
convention is what a bare matrix needs, because a bare matrix has no path to
follow at all.

\subsection{Contraction of charged tensors}
\label{sec:flow-contract}

Contracting two tensors along a shared pair of legs is the invariant pairing
induced by the CG data.  Two things must hold for the result to stay inside
the block-sparse format:
\begin{itemize}
  \item the shared legs must be oriented oppositely, so that the pairing is
        the invariant form $D$ (section~\ref{sec:conj}) and not the naive
        index sum;
  \item the total charge of the result must be a single irrep, i.e.\ the
        contraction must not mix several total charges.
\end{itemize}
The second condition is a property of the operands, not of the algorithm.  A
contraction of two invariant tensors is again invariant, so the matrix-product
operations - which contract the physical leg of a state with the physical leg
of an operator, and then the two virtual legs of the carry - never mix totals,
and the condition holds by construction.  It fails only for an object that is
not an intertwiner of a single total, such as a bare matrix $M_{ab}$ whose two
legs fuse to the trivial channel and to the adjoint channel at the same time;
such an object has to be carried in the resolved convention, where the two
channels are simply two components of one dense index pair.

The price of the resolved convention appears at the next step.  Because a
resolved block stores the magnetic index, it is the \emph{sum over all paths}
that reach the block''s leg sectors.  That is exactly what one wants for the
dense expansion of a network, and exactly what has to be undone to continue in
the sparse convention.  For a multiplicity-free group the sum has a single
term - the tree is determined by the leg sectors - and nothing is lost; for a
group with fusion multiplicities, e.g.\ the $\SU{3}$ adjoint product
$\mathbf 8\otimes\mathbf 8$ that contains $2\cdot\mathbf 8$, a block merges
several trees and the paths can no longer be read from the labels.

\subsection{The projection that recovers a fusion path}
\label{sec:flow-project}

A tensor with legs $\lambda_1,\dots,\lambda_k$ and total charge $c$ is expanded
in the CG data of its admissible fusion trees as
\begin{equation}
  T_{a_1\dots a_k} \;=\; \sum_{\text{trees }\tau}
  C^{\tau}_{a_1\dots a_k}\, T_{\tau},
  \label{eq:flow-expand}
\end{equation}
where the sum runs over the trees $\tau$ of the leg sectors that fuse to $c$
(including the copy index of every internal vertex when the group has fusion
multiplicities), $C^\tau$ is the vertex tensor of the tree, obtained by
multiplying the CG tensors along its chain, and $T_\tau$ is the reduced datum
of that tree.  This is Wigner--Eckart~\eqref{eq:we} written for a whole
network: the CG structure is not stored but looked up, and the block shape of
the reduced container is the product of the leg multiplicities alone.

The inverse of~\eqref{eq:flow-expand} is the projection.  Vertex tensors of
distinct trees are orthogonal,
\begin{equation}
  \sum_{a_1\dots a_k} \overline{C^{\tau}_{a_1\dots a_k}}\,
  C^{\tau'}_{a_1\dots a_k}
  \;=\; \lVert C^{\tau}\rVert^{2}\,\delta_{\tau\tau'},
  \qquad \lVert C^{\tau}\rVert^{2} = \dim V_{c},
  \label{eq:flow-orth}
\end{equation}
with the normalization~\eqref{eq:cg-norm} of the CG tensors, so that
\begin{equation}
  T_{\tau} \;=\; \frac{1}{\lVert C^{\tau}\rVert^{2}}
  \sum_{a_1\dots a_k} \overline{C^{\tau}_{a_1\dots a_k}}\, T_{a_1\dots a_k} .
  \label{eq:flow-project}
\end{equation}
Three details matter in an implementation.  First, the sum runs over the
magnetic index of every leg \emph{and} over the multiplicity index of every
leg: a leg that contains the same irrep $m$ times has $m$ independent copies,
each with its own dense window, and each copy carries one entry per tree.
Second, the total charge has to be a single irrep, for the reason of the
previous subsection: a block that sums several totals has no single
$\lVert C^\tau\rVert^{2}$ to divide by, and projecting it would keep one
channel and drop the others silently.  Third, the projection is the change of
basis between the two conventions, so it is exact up to rounding, and the
checks of this package verify it as a round trip of the dense expansion with
$\max|\Delta|\lesssim10^{-16}$.

\subsection{Splitting a charged tensor}
\label{sec:flow-split}

The other half of an application is the split: a tensor whose legs are (left
group, right group) is cut into a left factor and a right factor by a singular
value decomposition, and the bond between them carries the cut charge.  The cut
charge of a tree is the coupling of the left group inside the tree,
$c_{\text{cut}}=c_{|L|}$ in the notation of equation~\eqref{eq:tree}; in the
reduced convention it is read off the block label, in the resolved convention
from the leg sectors whenever the group is multiplicity-free.  Blocks are
grouped by $c_{\text{cut}}$, one matrix is assembled per group, and each is
decomposed separately,
\begin{equation}
  M_{q} \;=\; U_q\, S_q\, V_q^{\dagger},
\end{equation}
so that truncation happens per charge and the symmetry of the state survives
the truncation itself.

The bookkeeping of the new bond is where the resolved convention needs one
extra step.  If $r_q$ singular values are kept in the cut charge $q$, the new
bond holds $r_q$ copies of the irrep $q$: $r_q\dim V_q$ dense indices in the
resolved convention, $r_q$ entries in the reduced one.  The singular values
are data per copy, so each of them has to be repeated over the magnetic
components of its multiplet before it is folded into the right factor,
\begin{equation}
  s_{q}\,(V_q^{\dagger})_{\text{index}} \quad\text{for every index in the
  multiplet of } q ,
  \label{eq:flow-absorb}
\end{equation}
whereas in the reduced convention the factor''s index \emph{is} the copy and
the plain product is already correct.  Broadcasting $r_q$ numbers over
$r_q\dim V_q$ indices is a different operation, and it is right only when
$\dim V_q=1$ - the case of every Abelian group, and of the trivial irrep of
any group.  This is the step that a naive $s\,\text{reshape}(-1,1,1)$ gets
wrong, and it is why the split is a pair of operations, the decomposition and
the absorption of the singular values, rather than one.

\subsection{The three steps of the implementation}

The mechanism above is realized in three places.

\begin{enumerate}
  \item \emph{The projection.}  The engine walks the trees of the reduced
        labelling and, for every tree and every combination of leg
        multiplicity indices, writes the amplitude~\eqref{eq:flow-project}
        with the vertex tensor taken from the cached CG data of that tree.
        It belongs in the engine because the CG data and the tree enumeration
        live there; it cannot be done from the block keys of a resolved
        tensor in the Python layer, because those keys hold no internal
        vertices, so the paths to project onto are not in the key.  The Python
        layer only decides when a tensor is resolved.
  \item \emph{The absorption.}  The second factor of a split takes each kept
        singular value repeated over the multiplet of its cut charge, as
        in~\eqref{eq:flow-absorb}; the two callers are the two-site gate and
        the operator application.  The reduced route needs no repetition, so
        the operation is a no-op there.
  \item \emph{The refusal.}  A charge-changing operator is now accepted on any
        state, in either convention, because the projection makes the reduced
        route exact - for the multiplicity-free groups and for $\SU{3}$ with a
        bond of multiplicity two alike
        (section~\ref{sec:anchors}).  What stays refused is what the container
        genuinely cannot represent: a one-site operator whose charge is
        non-trivial and which therefore has no path to follow, and a reduced
        tensor whose content mixes several total charges.
\end{enumerate}

\subsection{Abelian groups, point groups, and the general case}
\label{sec:flow-groups}

Nothing in the mechanism is specific to $\SU{2}$ or $\SU{3}$: it uses the
fusion rules, the CG (vertex) data with its normalization, and the invariant
form that orients a leg.  What changes from group to group is which of the
three degenerates.

\paragraph{Abelian groups.}  Every irrep is one-dimensional, so the sector
dimension is one, the CG data are trivial (every coefficient is $1$), the
fusion is multiplicity-free and the invariant form is the identity.  The
resolved and the reduced convention then coincide index by index: the
projection~\eqref{eq:flow-project} is the selection of the block that carries
the charge, the repetition~\eqref{eq:flow-absorb} is the identity, and a
charge-changing operator is exact in either route.  This is the case of a
$U(1)$ charge (particle number, $S^z$), of $\mathbb Z_n$ and of the fermion
parity, and it is why the sparse machinery was exact for them long before the
non-Abelian case was.

\paragraph{Point groups.}  A crystallographic point group is finite: finitely
many irreps, of dimension $1$, $2$ or $3$.  The fusion rules are the
Clebsch--Gordan series of the group, computable from the character table, and
the CG data are the orthogonal matrices that reduce
$V_\lambda\otimes V_\mu$.  Two features of the general case appear here: a
multiplicity $N^{\nu}_{\lambda\mu}>1$ whenever the character sum
$\frac{1}{|G|}\sum_g\chi_\lambda(g)\chi_\mu(g)\overline{\chi_\nu(g)}$ exceeds
one, and a Frobenius--Schur indicator that separates the orthogonal from the
symplectic irreps (and, for a double group, the spinorial ones).  Both are
handled by the same three ingredients, with $\dim V_\lambda$ read from the
character table.

\paragraph{The general case.}  Any symmetry whose irreps carry a fusion ring
with a CG calculus - a compact Lie group, a finite group, a quantum group at a
root of unity - is covered by the same three ingredients, because neither the
projection nor the split ever asks what the group is: they ask for the sectors
of a leg, the trees between them, the vertex tensor of a tree, and the
invariant form.  In this package that is exactly the interface of the group
object - its fusion outcomes, its fusion tensor, its sector dimensions, its
invariant form and its duality matrix - and the engine branches on none of
them.  Adding a group therefore means adding its fusion and CG data, not a
case in the algorithm.  The design rule is that the symmetry logic is
universal and sealed inside the group layer, and the charge flow described
here is the sharpest test of it, because it is the operation that needs every
piece of the group data at once.


\section{Where the dense route belongs}
\label{sec:reference}

The package keeps a dense route: tensors whose legs store the magnetic index
explicitly (the ``resolved'' convention) hold ordinary dense blocks, and
contract by ordinary dense matrix multiplication.  It serves three purposes:

\begin{enumerate}
  \item \emph{An oracle.}  Any sparse network can be expanded and compared
        against the dense result; the checks of this package all do this,
        which is how convention errors of the kind in
        section~\ref{sec:conj} are caught.
  \item \emph{A bridge.}  A user''s bare matrix, a string operator, or an
        operator given only as a dense array can be turned into a symmetric
        object in the resolved convention and then manipulated.
  \item \emph{A slow but general fallback.}  For a group whose fusion data
        are not implemented, or for a probe of a small tensor, the dense
        route always works.
\end{enumerate}

It is not the production path.  A dense-block contraction restores the full
$D^2$ cost per leg and materializes the CG structure that the sparse route
never stores.  The intended split is therefore:

\begin{center}
\begin{tabular}{ll}
  \textbf{production} & block-sparse labels: charges on the legs, blocks
  labelled by trees, CG inserted only when a leg is expanded \\
  \textbf{reference} & resolved legs with dense blocks: plain dense
  arithmetic, used by tests, conversions and small diagnostics \\
\end{tabular}
\end{center}

\section{Measured anchor points}
\label{sec:anchors}

The following numbers were measured in this package and are quoted here as
the empirical anchor of the statements above.

\begin{itemize}
  \item $\SU{2}$ CG normalization~\eqref{eq:cg-norm}: with a unit placement the
        tree realizes $C_s/\dim V_s$ exactly, i.e.\ $\lVert C_s\rVert^2 =
        \dim V_s = 1,3,5$ for the sectors $(0),(2),(4)$.
  \item $\SU{3}$ dimension and Casimir: the constructed irreps reproduce
        $\dim V_{(p,q)}=(p+1)(q+1)(p+q+2)/2$ and
        $C_2=(p^2+q^2+pq+3p+3q)/3$ for every $(p,q)$ with $p+q\le4$, and the
        Klimyk multiplicities~\eqref{eq:klimyk} reproduce
        $\mathbf 8\otimes\mathbf 8=\mathbf{27}\oplus\mathbf{10}\oplus
        \overline{\mathbf{10}}\oplus2\cdot\mathbf 8\oplus\mathbf 1$.
  \item Frobenius--Schur~\eqref{eq:fs}: putting the sign in the wrong place
        (the free-function convention instead of the group''s duality
        matrix) turns the effective Hamiltonian non-Hermitian; the
        measurement that located it was a negative norm of an environment,
        while every symmetry-blind check stayed green.
  \item Sparsity: the same $\SU{2}$-symmetric state at $L=16$ with bond
        dimension $64$ needs $3806$ numbers against $43688$ for the
        symmetry-blind representation, and its sweep is $8.5$ times faster.
  \item The dense route: a resolved tensor built from a bare matrix
        round-trips exactly, $T\!\to\!\text{dense}\!\to\!T$ with
        $\max|\Delta|=0$, for $S^z$, $S^x$, the identity of $\SU{2}$ and
        $\lambda_8$ of $\SU{3}$ - which is what makes it usable as an oracle.
  \item Charge flow, $\SU{2}$: applying an operator of total charge $1$ (a
        triplet, the sector $(2)$) to the $L=4$ Heisenberg ground state gives
        $\lVert\mathcal O\psi\rVert^2=0.5290063509461093$ against
        $0.5290063509461105$ from the dense expansion of the same operator, and
        at $L=6$ with two such operators $1.1901684550722649$ against
        $1.1901684550722691$.  The bonds of the result carry
        $(1)$, $(0)\oplus(2)$, $(1)\oplus(3)$,
        $(0)\oplus(2)\oplus(4)$, $(1)$: the charge the operator added, and not
        the trivial one.
  \item Charge flow with a multiplicity, $\SU{3}$: the same application on the
        $L=6$ ULS chain, whose bond between sites $2$ and $3$ holds two copies
        of the adjoint, gives $0.10983847951600069$ against
        $0.10983847951600301$ dense for the adjoint channel and
        $0.0888259974999079$ against $0.08882599749990971$ for the
        $\mathbf{10}$ channel, so the multiplicity is carried exactly instead
        of being dropped.
  \item The projection~\eqref{eq:flow-project}: applied to a merged state and
        operator tensor it reproduces the reduced labelling of the same
        contraction to $1.2\times10^{-16}$, with the same block keys, and both
        contractions agree with the plain dense contraction of the three
        tensors to $3\times10^{-16}$.
  \item The absorption~\eqref{eq:flow-absorb}: without the repetition of each
        singular value over its multiplet, the application fails with a size
        mismatch ($4$ against $2$) as soon as a kept charge has
        $\dim V_q>1$; with the repetition, and with the projection above, a
        charge-changing operator applies in either convention and the checks
        of the package are green.
  \item Where the time goes: one $\SU{2}$ sweep at $L=10$ with bond dimension
        $32$ issues $1.8\times10^{5}$ CUDA kernels of $2.3\,\mu$s on average,
        $2.3\times10^{4}$ device-to-device copies and $1.6\times10^{5}$ block
        slices, for a total device time of $0.21$ s against a wall time of
        $1.4$ s.  The GPU is therefore idle for most of the sweep: the cost is
        the dispatch of one tiny operation per block, not the arithmetic.  The
        decomposition is negligible in that budget ($36$ SVD calls, $0.7\%$ of
        the CPU time), and the default intra-op thread pool does not pay for
        itself - the same sweep takes $2.9$ s with the pool and $1.4$ s
        without it, because a workload of tiny operations cannot be spread
        over threads.
\end{itemize}
